\subsection{W1包络结果与既有比较}
\label{sec:w1-enclosure-results}
两种规则分别在31533和37671个活动单元达到全部七函数的预设精度；按协议预留口径合计1351046次评价，未新增MCMC拟合。完整原件的24份资产、两个区域的覆盖、叶单元求积与余项累计、有符号比率和交集均通过精确有理端点核对。两种求积的共同实现范围见第\ref{sec:w1-enclosure-method}节。

\begin{table}[htbp]\centering\small
\begin{tabular}{lrl}\toprule
函数 & 后验期望区间 & Gauss/Simpson达标\\\midrule
$\alpha$ & $[0.606577102075, 0.606577121451]$ & 是/是\\
$\beta$ & $[-0.622983259465, -0.622983239584]$ & 是/是\\
$\alpha^2$ & $[0.371577048819, 0.371577061368]$ & 是/是\\
$\beta^2$ & $[0.397618700920, 0.397618714571]$ & 是/是\\
$p(0)$ & $[0.647037387034, 0.647037407034]$ & 是/是\\
$p(100)$ & $[0.495902573466, 0.495902588876]$ & 是/是\\
$\Pr(\beta>0)$ & $[0.000000000047671124, 0.000000000047673026]$ & 是/是\\
\bottomrule\end{tabular}
\caption{W1后验函数的数值包络交集。端点向外舍入；各规则的精度状态由未舍入区间判断。连续函数半宽目标为$10^{-8}$，事件半宽目标为$10^{-13}$且相对半宽不超过1\%。}
\label{tab:w1-enclosure}\end{table}
两种规则的七函数均达到预设目标。


表中连续函数以12位小数、事件以18位小数向外舍入；精度是否达标使用原二进制区间判断，因此十进制展示宽度可能略大于停止阈值。原七个求积参考点全部位于相应交集内，原参考文件和主结果表保持不变。

对原432份有效W1拟合（九工作流、两预算，每格24份四链重复），传播区间得到126项单方法平方差范围和504项配对损失差范围。连续六函数各有68项符号固定、4项恒等；事件的72项比较全部恒等。合计408项符号固定、96项恒等，没有相对原参考点的符号翻转或跨零未定项。数值符号稳定包含极小的浮点路径差异，并不说明效应有实际意义或在重复实验中显著；主实验的逐点不确定性区间继续单列。

全部W1事件估计为零，而表中概率区间严格为正，故事件估计相对该参考的误差为100\%。其绝对平方差约为$2.27\times10^{-21}$，说明极小的绝对损失仍可与未观察到罕见事件同时出现；常量函数的链间诊断保持未定。

恢复检查复用全部24份原科学资产，新增被积函数评价为零，文件哈希未变。本地归档含81个内容成员；在同机独立目录从归档源码和输入重建后，九份核验/报告文件及两份校验清单逐字节相同。该检查确认可搬移重建，不替代对导数与尾界实现的数学条件，也不是异机独立研究复现。
